\subsection{backup.var}
\texttt{backup.var} configures what gets preserved/uploaded at the end of a run and where.

\subsubsection{Upload / backup switches}
\begin{lstlisting}
UPLOAD=0/1     : upload figures/ASCII outputs to the remote server(s). [current: 1]
ARUPLOAD=0/1   : rsync the AUTOREGRESSION hand-off directory to the short-timescale
                 component's host(s) at the end of run.sh, if that component is in use
                 on this deployment. [current: 1]
BACKUP=0/1     : perform the local backup of output files. [current: 1]
BACKUPFIGURES=0/1: keep a local mirror of the delivered figures. [current: 1]
\end{lstlisting}

\subsubsection{Backup parameters}
\begin{lstlisting}
TAIL_DIR="string"   : suffix appended to the per-night backup directory name. The
                      checked-in value is a short project tag inherited from this
                      codebase's origin; edit to identify your own deployment.
BACKUP_DIR="string" : root of the local ASCII/plot backup. [current: "${HOME}/BACKUP_OUTPUT"]
BACKUP_FIGURES_DIR="string": root of the local figures mirror.
                      [current: "${HOME}/BACKUP_FIGURES"]
BACKUP_REAL, BACKUP_RUN,
BACKUP_DIAG=0/1     : whether to keep, respectively, the PREP_REAL / EXE(RUN) / DIAG
                      binary output in the per-night backup (rather than discarding it).
                      [current: 0, 0, 0]
BACKUP_MNHDAT=0/1   : keep the ASCII (.dat, FICVAL, AUTOREGRESSION) postprocessing
                      output. [current: 1]
BACKUP_POSTPROC=0/1 : keep the ASCII+PLOT postprocessing output tree. [current: 1]
BACKUP_LOG=0/1      : keep the run's log directory (required for backup.sh to be able
                      to relocate its own log paths correctly at the end of the run).
                      [current: 1]
\end{lstlisting}
The corresponding \texttt{BACKUP\_*\_NAME} variables are the actual destination paths, all built from
\texttt{BACKUP\_BASE\_NAME} (itself \texttt{BACKUP\_DIR/<date><TAIL\_DIR>\_<N>DOM\_<MESONH version>});
they are not normally edited directly.

\subsubsection{Remote host parameters}
\begin{lstlisting}
DAYSTARTBEFORE=integer (negative): how many days before the forecast date the simulation
                      started; used to locate the right backup sub-directory.
                      [current: -1]
REMOTE_USER="string", REMOTE_HOST="string": account/host of the primary remote output
                      server. Placeholder values in the checked-in file -- must be
                      filled in for your own deployment.
REMOTE_USER_BK="string", REMOTE_HOST_BK="string": account/host of the redundant backup
                      remote server. Also placeholders in the checked-in file.
REMOTE_PATH_ARCHIVE="string", REMOTE_PATH_OUTPUTS_GROUND="string": upload paths on the
                      primary server for figures and ground-parameter ASCII output
                      respectively (trailing slash required). The checked-in values
                      still point at a path tree named after this codebase's origin;
                      edit to match your own server layout.
REMOTE_PATH_ARCHIVE_BK="string",
REMOTE_PATH_OUTPUTS_GROUND_BK="string": matching upload paths on the backup server.
REMOTE_AR="string", REMOTE_AR_BK="string": address(es) of the short-timescale
                      component's hand-off host(s), if that component is deployed;
                      the AUTOREGRESSION directory is rsynced there at the end of
                      run.sh. Placeholders in the checked-in file.
REMOTE_PATH_AUTOREG="string": remote path (trailing slash) for that hand-off rsync.
\end{lstlisting}
